\scanpage{89}

\uline{1.4 suite}

Choisissons d'abord un splitting de \inserted{$\mathscr L$ i.e. de} $\unsure{\mathfrak g|\mathcal G}$, un

\[ \mathfrak g\simeq\Omega+\underset{\substack{\Vert\\\underline O_X}}{\deleted{\unsure{O}}\Omega^0}+\Omega^{-1} \]

d'où, posant $\omega=\underline\Omega^1_{X/S}$
\[ \mathfrak g\otimes\omega\simeq\Omega\otimes\omega+\omega+\Omega^{-1}\otimes\omega \]

\deleted{\unsure{On en a une}} Choisissons-en loc. une section $\deleted{\unsure{\omega}}\delta$ de $\Omega$ comme base. Une connexion (sans courbure) $c_0$ est obtenue par

\[ \begin{aligned}
 \mathfrak g&\xrightarrow{c_0}\mathfrak g\otimes\omega\\
 c_0(\delta)&=c_0(1)=c_0(\delta^{-1})=0
 \end{aligned} \]

i.e.
\[ c_0\bigl(\underbrace{\lambda\delta+\mu.1+\nu\delta^{-1}}_{\substack{\updownarrow\\(\lambda,\mu,\nu)}}\bigr)=\underset{\substack{\updownarrow\\(d\lambda,d\mu,d\nu)}}{\delta\,d\lambda+1\,d\mu+\delta^{-1}d\nu} \]

Une connexion générale s'obtient alors en \unsure{ajoutant|rajoutant} une section de
\[ \mathrm{Dér}(\mathfrak g,\mathfrak g)\otimes\omega\simeq\mathfrak g\otimes\omega \]
(\uline{NB} toute dérivation de $\mathfrak g$ est intérieure !)

Soit \inserted{$\varpi_{\unsure{*|\bullet}}={}$} $\varpi_2\oplus\varpi_1\oplus\varpi_0\in\Gamma\mathfrak g\otimes\omega$ i.e.
\[ \begin{aligned}
 \varpi_2&\in\Gamma(\Omega\otimes\omega)\\
 \varpi_1&\in\Gamma(\omega)\\
 \varpi_0&\in\Gamma(\Omega^{-1}\otimes\omega)
 \end{aligned} \]

on aura donc \unsure{pour|par} connexion $c^{\deleted{\unsure{\varpi}}}_{\varpi_*}$ définie par $c_0$ et $\varpi_*$ 
\[ c_{\varpi_*}\bigl(\underbrace{\lambda\delta+\mu.1+\nu\delta^{-1}}_{(\lambda,\mu,\nu)}\bigr)\overset{?}{=}\deleted{\unsure{\delta\,d\lambda+1\,d\mu+\delta^{-1}d\nu}} \]

\[ \unsure{+|\ast} \]

% archive sheet number: 88
